\chapter{Abordări anterioare}
\section{Rețele convoluționale pentru clasificarea sunetelor}
Prin experimentele celor de la Google prezentate în \cite{large} s-a demonstrat că rețelele de mari dimensiuni
ce au la bază straturi convoluționale precum ResNet-50 și AlexNet au rezultate mai bune în comparație 
cu rețelele cu straturi dens conectate și au fost utilizate pentru clasificarea a până la 30 de mii de 
sunete prezente într-un set de date de 70 de milioane de videoclipuri.
\section{Clasificarea genurilor muzicale}
Odată cu creșterea popularității platformelor de streaming audio, problema clasificării genurilor muzicale a devenit 
,de asemenea, mai populară. 

În \cite{mgcrelated} au fost comparate performanțele rețelelor convoluționale ce primesc ca input 
spectrograme, cu cele ale algoritmilor SVM care primesc ca input caracteristici extrase manual din semnalul audio și apoi s-a incercat 
combinarea acestora.
Experimentele s-au bazat pe trei seturi de date diferite, rezultatele obținute pe două dintre acestea pentru combinația dintre CNN și SVM
depășind rezultatele existente precendent în literatură.

O altă abordare des întâlnită este combinarea rețelelor convoluționale cu rețele neuronale recurente. Experimentele au demonstrat că un CNN 
combinat cu o rețea recurentă are rezultate mai bune decât rețetele convoluționale simple \cite{resut1} \cite{resut2}.

În ceea ce privește clasificarea genurilor muzicale pe baza elementelor vocale, cele mai multe experimente s-au axat pe clasificarea artiștilor, 
nefiind cazul de o separare a surselor sonore  pentru a izola vocea. Un experiment care se concentreză exclusiv pe elementele vocale este 
prezentat în \cite{vocal}. Cu toate acestea, utilizează caracteristici de dimensiuni mici extrase din semnalul audio și un algoritm de 
clasificare SVM, iar experimentele prezentate anterior demonstrează faptul că rezultate mai bune pot fi obținute cu rețele neuronale și 
reprezentări de dimensiuni mai mari ale semnalelor audio precum spectograme sau coeficienti MFCC.